\documentclass[11pt, a4paper]{article}

% ====== Required Packages ======
\usepackage[utf8]{inputenc}
\usepackage[T1]{fontenc}
\usepackage{lmodern}
\usepackage{microtype}
\usepackage{amsmath, amssymb}
\usepackage{geometry}
\usepackage{booktabs}
\usepackage{tabularx}
\usepackage{array}
\usepackage[table]{xcolor}
\usepackage{hyperref}
\usepackage{enumitem}
\usepackage{fancyhdr}
\usepackage{listings}
\usepackage{titlesec}

\geometry{top=2.4cm, bottom=2.4cm, left=2.4cm, right=2.4cm, headheight=14pt}

% ====== Colors ======
\definecolor{accentblue}{RGB}{30, 80, 160}
\definecolor{passgreen}{RGB}{198, 239, 206}
\definecolor{warnyellow}{RGB}{255, 235, 156}
\definecolor{warnred}{RGB}{255, 199, 206}
\definecolor{tableheader}{RGB}{230, 236, 245}
\definecolor{codebg}{RGB}{246, 248, 250}
\definecolor{codekw}{RGB}{170, 13, 145}
\definecolor{codecmt}{RGB}{0, 110, 60}
\definecolor{codestr}{RGB}{196, 26, 22}

\titleformat{\section}{\large\bfseries\color{accentblue}}{\thesection}{1em}{}
\titleformat{\subsection}{\normalsize\bfseries}{\thesubsection}{1em}{}

\hypersetup{
  colorlinks=true,
  linkcolor=accentblue,
  urlcolor=accentblue,
  pdftitle={DAMPS-MMHCL Wave 1 / M1.5 -- Sweep-Cell Patch and Post-hoc Analysis Cell (Instructions)},
}

\lstdefinestyle{py}{
  language=Python,
  backgroundcolor=\color{codebg},
  basicstyle=\ttfamily\scriptsize,
  keywordstyle=\color{codekw}\bfseries,
  commentstyle=\color{codecmt}\itshape,
  stringstyle=\color{codestr},
  numbers=left,
  numberstyle=\tiny\color{gray},
  numbersep=8pt,
  showstringspaces=false,
  breaklines=true,
  frame=single,
  rulecolor=\color{gray!40},
  tabsize=4,
  columns=fullflexible,
  literate={->}{{$\rightarrow$}}1 {>=}{{$\geq$}}1 {<=}{{$\leq$}}1
}

\newcommand{\cmark}{\textcolor{green!60!black}{\textbf{\checkmark}}}
\newcommand{\warnmark}{\textcolor{orange!80!black}{\textbf{!}}}

\newcolumntype{Y}{>{\centering\arraybackslash}X}
\newcolumntype{Z}{>{\raggedright\arraybackslash}X}

\pagestyle{fancy}
\fancyhf{}
\fancyhead[L]{\small\textit{DAMPS-MMHCL Wave 1 / M1.5 --- Patch Instructions}}
\fancyhead[R]{\small\thepage}
\fancyfoot[C]{\small Phase 2 Execution Hygiene Report}
\renewcommand{\headrulewidth}{0.4pt}
\renewcommand{\footrulewidth}{0.4pt}

\title{%
  \textbf{Wave 1 / M1.5 Sweep --- Notebook Patch \& Post-hoc Analysis Cell}\\[0.4em]
  \large DAMPS-MMHCL Phase 2 Section 8.2 --- Execution Hygiene Instructions\\[0.2em]
  \normalsize Replace the final notebook cell and append a post-hoc analysis cell
}
\author{DAMPS-MMHCL Research Team}
\date{\today}

\begin{document}
\maketitle
\tableofcontents
\newpage

% ============================================================
\section{Motivation}
% ============================================================
While inspecting the most recent M1.5 sweep logs we found two execution-hygiene
defects that are far more serious than the early-stopping question they were
originally raised under:

\begin{enumerate}[leftmargin=*,itemsep=4pt]
  \item \textbf{The runs are not a real M1.5 LogQ sweep.}
        The attached notebook (\texttt{file:103}) contains \emph{zero}
        occurrences of \texttt{enable\_logq}, \texttt{logq\_scale}, or
        \texttt{M1.5}. The launched command lines still consist solely of the
        nine \texttt{-{}-damps\_*} Phase~1 flags. Worse, the log timestamps read
        \texttt{2026-05-12}, i.e.\ they are recycled May outputs --- not fresh
        runs with the LogQ correction. The
        $\text{Recall@}20 \approx 0.088\text{--}0.091$ numbers therefore reflect
        the bare rev45 baseline; no LogQ delta has been applied yet.
  \item \textbf{Seeds are corrupted by an implicit \texttt{float} cast.}
        The logs contain \texttt{seed=1.404632241e+09}, \texttt{seed=1.569402867e+09},
        i.e.\ the seed was coerced to scientific-notation \texttt{float} before
        reaching \texttt{-{}-seed}. This silently round-trips through 32-bit
        rounding and breaks bit-level reproducibility of the sweep. The fix is
        \texttt{str(int(seed))} at the call site.
\end{enumerate}

The early-stopping logic itself is sound. If you want runs to terminate more
cleanly after convergence (to save GPU time), you may lower
\texttt{early\_stopping\_patience} from \textbf{30} to \textbf{20}; this is
optional and not required for correctness.

% ============================================================
\section{Acceptance Criteria for the Patched M1.5 Sweep}
% ============================================================
\begin{itemize}[leftmargin=*,itemsep=3pt]
  \item Every launched \texttt{train.py} command must contain all five LogQ
        flags, namely:
\begin{lstlisting}[style=py,language=bash,numbers=none]
--enable_logq 1
--logq_mode {laplace, raw, sqrt}
--logq_beta  1.0
--logq_scale {0.05, 0.1, 0.3, 1.0}
--logq_clip  5.0
\end{lstlisting}
  \item Every \texttt{-{}-seed} argument must be the string form of an integer
        (e.g.\ \texttt{1404632241}), never scientific notation.
  \item The M1.5 gate must be evaluated on \texttt{BEST\_Test\_Recall@20}
        (the test-set snapshot taken at the validation-recall peak epoch),
        \emph{not} on any mid-run epoch printout.
  \item The per-run logger directory --- whose name embeds
        \texttt{seed} and \texttt{ablation\_target} --- must encode
        \texttt{(mode, scale, seed)} so configurations do not overwrite each
        other's log files.
  \item Validation-recall curves for every $(\text{mode}, \text{scale},
        \text{seed})$ must be plotted post-hoc so it can be visually verified
        that no run was terminated within $<15$ epochs of its recall peak.
\end{itemize}

% ============================================================
\section{Part 1 -- Replace the Final Notebook Cell (M1.5 Sweep Driver)}
% ============================================================
Open the attached notebook and overwrite the contents of the last code cell
(the cell titled ``Section~8.2 -- Wave~1 / M1.5 sweep driver'') with the
script reproduced verbatim in Appendix~A
(\texttt{cell30\_m15\_fixed.py}). The patched cell embodies five concrete
changes relative to the previous revision:

\begin{table}[h]
\centering
\renewcommand{\arraystretch}{1.25}
\rowcolors{2}{white}{tableheader!35}
\begin{tabularx}{\textwidth}{@{}p{0.36\textwidth}Z@{}}
\toprule
\rowcolor{tableheader}
\textbf{Defect in previous cell} & \textbf{Fix applied in the patched cell} \\
\midrule
LogQ flags missing -- launched command line is rev45 Phase~1 only &
The constant \texttt{BACKBONE\_FLAGS} now omits LogQ; LogQ activation
(\texttt{-{}-enable\_logq 1}, \texttt{-{}-logq\_mode}, \texttt{-{}-logq\_beta},
\texttt{-{}-logq\_scale}, \texttt{-{}-logq\_clip}) is appended unconditionally
inside \texttt{run\_one(mode, scale, seed)}. \\
\texttt{seed} coerced to \texttt{float} ($\sim$\texttt{1.4e+09}) &
\texttt{\_as\_int\_seed = lambda s: int(round(float(s)))} is applied when the
seed list is bound, and \texttt{str(int(seed))} is used at the call site for
\texttt{-{}-seed}. \\
Sweep configs collide in the per-run log directory because
\texttt{ablation\_target} is the constant \texttt{wave1\_m15\_logq} &
\texttt{ablation\_target = f"m15\_logq\_\{mode\}\_s\{scale\}\_seed\{seed\}"}.
This propagates into \texttt{train.py}'s \texttt{\_experiment\_paths()} and
yields a unique directory per configuration. \\
M1.5 gate previously read mid-run epoch metrics & Regex parser now anchors on
the \texttt{BEST\_Test\_*} family
(\texttt{BEST\_Test\_Recall@20}, \texttt{BEST\_Test\_NDCG@20},
\texttt{BEST\_Test\_Precision@20}) and additionally extracts
\texttt{BEST\_Val\_Recall\_Peak\_Epoch} /
\texttt{BEST\_Val\_NDCG\_Peak\_Epoch} for the post-hoc audit. \\
Runs over-train after convergence & \texttt{early\_stopping\_patience} default
lowered from \textbf{30} to \textbf{20}; flip \texttt{BACK\_TO\_PATIENCE\_30
= True} to reproduce the previous behaviour. \\
\bottomrule
\end{tabularx}
\caption{Five defects in the previous M1.5 sweep cell and their corresponding
fixes in the patched cell.}
\end{table}

Verbatim hyper-parameters (from \texttt{Phase 2 Section 8.2}) used by the
patched cell:

\begin{itemize}[leftmargin=*,itemsep=2pt]
  \item \texttt{LOGQ\_SCALES = [0.05, 0.1, 0.3, 1.0]} (four-point scale sweep);
        \texttt{LOGQ\_MODES = ["laplace"]} (extend to \texttt{["raw","sqrt"]}
        for the full 24-run sweep).
  \item \texttt{N\_SEEDS = 5}; \texttt{LOGQ\_BETA = 1.0}; \texttt{LOGQ\_CLIP = 5.0}.
  \item Acceptance gate: $\text{mean BEST\_Test\_R@20} \ge 0.0925$ \textbf{and}
        $\min \text{BEST\_Test\_R@20} \ge 0.0890$.
\end{itemize}

\subsection*{Reporting protocol}
When you summarise an M1.5 outcome, quote exclusively
\texttt{BEST\_Test\_Recall@20} (the val-peak snapshot). Do \emph{not} cite
$\text{Recall@}20$ printed on epoch 84 or 219 in the middle of a run --
those numbers are explicitly \emph{not} the model that would be deployed under
the \texttt{restore\_best=1} early-stopping policy.

% ============================================================
\section{Part 2 -- Append a Post-hoc Analysis Cell}
% ============================================================
Append the script in Appendix~B (\texttt{cell31\_m15\_analysis.py}) as the
\textbf{new last cell} of the notebook (preceded by a level-2 Markdown header
titled ``Wave~1 / M1.5 Sweep -- Post-hoc Analysis''). The cell is read-only:
it never launches \texttt{train.py}; it only inspects log artefacts already
produced by the sweep cell.

\subsection*{What it does}
\begin{enumerate}[leftmargin=*,itemsep=3pt]
  \item Globs every directory under \texttt{../<dataset>/damps\_*} whose name
        matches the regex below, and locates its \texttt{.txt} logger output:
\begin{lstlisting}[style=py,language=bash,numbers=none]
damps_.._seed=<seed>_m15_logq_<mode>_s<scale>_seed<seed>
\end{lstlisting}
  \item Parses each log for the trailing summary lines emitted by
        \texttt{train.py} at the end of training, namely:
\begin{lstlisting}[style=py,language=bash,numbers=none]
BEST_Test_Recall@20 / NDCG@20 / Precision@20
BEST_Val_Recall@20  / NDCG@20
BEST_Val_Recall_Peak_Epoch
BEST_Val_NDCG_Peak_Epoch
\end{lstlisting}
        plus every per-epoch diagnostic line of the form:
\begin{lstlisting}[style=py,language=bash,numbers=none]
Epoch N [..s + ..s]: loss=.. recall@10=.. recall@20=.. ndcg@20=..
\end{lstlisting}
  \item Builds a per-run audit table containing, for every
        $(\text{mode}, \text{scale}, \text{seed})$, the
        \texttt{BEST\_Test\_Recall@20}, the validation-recall peak epoch, the
        last epoch reached, and the margin
        \texttt{(last - peak)}. Any run whose margin is below 15 epochs is
        flagged for inspection (this would be evidence that early stopping
        cut the run short of its real peak).
  \item Re-computes the M1.5 gate purely from the parsed logs, as a sanity
        cross-check of the in-memory \texttt{all\_runs} dictionary built by the
        sweep cell.
  \item Plots, for every \texttt{logq\_scale}, the val/recall@20 curve of each
        seed on a shared axis; overlays a dotted vertical line at the recall
        peak, plus horizontal lines at the rollback gate ($0.0890$) and the
        M1.5 gate ($0.0925$). The figure is saved as
        \texttt{../<dataset>/m15\_logq\_val\_recall\_curves.png}.
  \item Persists the parsed run dictionaries and aggregate to
        \texttt{../<dataset>/m15\_logq\_parsed\_logs.json}.
\end{enumerate}

\subsection*{Interpreting the visual output}
For every (scale, seed) curve, look for two patterns:

\begin{itemize}[leftmargin=*,itemsep=3pt]
  \item \cmark\ \textbf{Healthy:} the curve rises monotonically (with
        validation noise), reaches its peak around the middle or back third of
        the run, then plateaus or drops; the last epoch is at least 15--20
        epochs past the peak. The dotted vertical line sits comfortably to the
        left of the right edge. \texttt{restore\_best=1} guarantees that the
        \emph{deployed} model is the one at the peak, not the final epoch.
  \item \warnmark\ \textbf{Suspect:} the peak is within the last 5--10 epochs
        of the run, i.e.\ the curve still has positive slope when training
        ended. This would indicate the run was halted prematurely. The cell
        flags such cases in the console output.
\end{itemize}

% ============================================================
\section{Part 3 -- Recipe to Execute the Patch}
% ============================================================
\begin{enumerate}[leftmargin=*,itemsep=3pt]
  \item Open the notebook
        \texttt{Local\_Random\_seeds\_train\_mmhcl\_clothing\_colab\_original.ipynb}
        in JupyterLab.
  \item Replace the contents of the final code cell with
        \texttt{cell30\_m15\_fixed.py} (Appendix~A).
  \item Append a new Markdown header (\verb|## 8. Wave 1 / M1.5 Sweep -- Post-hoc Analysis|)
        and a new code cell containing \texttt{cell31\_m15\_analysis.py}
        (Appendix~B).
  \item Save the notebook as
        \texttt{Local\_Random\_seeds\_train\_mmhcl\_clothing\_colab\_rev54\_wave1\_fixed.ipynb}.
  \item Re-run the sweep cell (the original took $\sim$20 runs $\times$ 4--6h
        on RTX 5090). Once it finishes, run the analysis cell on the same
        kernel.
  \item Quote \texttt{BEST\_Test\_Recall@20} (not mid-run printouts) when
        reporting the M1.5 gate verdict, and attach
        \texttt{m15\_logq\_val\_recall\_curves.png} as evidence that no run
        was terminated prematurely.
\end{enumerate}

% ============================================================
\appendix
\section{Appendix A -- Patched Cell (\texttt{cell30\_m15\_fixed.py})}
% ============================================================
\lstinputlisting[style=py]{cell30_m15_fixed.py}

\newpage
% ============================================================
\section{Appendix B -- Post-hoc Analysis Cell (\texttt{cell31\_m15\_analysis.py})}
% ============================================================
\lstinputlisting[style=py]{cell31_m15_analysis.py}

\end{document}
